\chapter{The Trust Problem}
\label{ch:trust-problem}

% Working draft - S. Vene, 2026-02-08

\section{Introduction: Why Trust Matters}
\label{sec:trust-matters}

The adoption of AI agents in enterprise environments is fundamentally constrained by trust. Unlike traditional automation, which executes deterministic rules predictably, AI agents exhibit autonomous decision-making that introduces uncertainty. Organizations face a dilemma:

\begin{itemize}
    \item \textbf{Without trust:} Agents are limited to trivial tasks, losing their value proposition
    \item \textbf{With misplaced trust:} Agents can cause significant damage through unexpected behavior
    \item \textbf{With calibrated trust:} Agents can be deployed effectively with appropriate safeguards
\end{itemize}

This chapter examines the trust problem in depth, proposing a theoretical framework for reasoning about trust in agentic systems.

\section{The Deterministic-Probabilistic Shift}
\label{sec:deterministic-probabilistic}

\subsection{Traditional Automation: Deterministic Systems}
\label{subsec:deterministic-systems}

Traditional CI/CD pipelines and infrastructure automation operate on deterministic principles:

\begin{center}
Input $\rightarrow$ Fixed Rules $\rightarrow$ Output
\end{center}

Given the same input and configuration, the system produces the same output. This determinism enables:
\begin{itemize}
    \item \textbf{Predictability:} Teams know what will happen
    \item \textbf{Testability:} Behavior can be verified before deployment
    \item \textbf{Debugging:} Failures can be traced to specific inputs/rules
    \item \textbf{Compliance:} Auditors can verify behavior through rule inspection
\end{itemize}

\subsection{Agentic Systems: Probabilistic Behavior}
\label{subsec:probabilistic-behavior}

AI agents fundamentally break this model:

\begin{center}
\begin{tabular}{c}
Input $\rightarrow$ LLM Reasoning $\rightarrow$ Probabilistic Output \\
$\uparrow$ \\
Context, Temperature, Model State
\end{tabular}
\end{center}

The same input can produce different outputs due to:
\begin{itemize}
    \item Stochastic sampling in LLM inference
    \item Context window variations
    \item Model updates and drift
    \item Emergent reasoning patterns
\end{itemize}

This shift from deterministic to probabilistic has profound implications for trust:

\begin{table}[htbp]
\centering
\caption{Deterministic vs Probabilistic Systems}
\label{tab:deterministic-probabilistic}
\begin{tabular}{lll}
\toprule
\textbf{Property} & \textbf{Deterministic} & \textbf{Probabilistic} \\
\midrule
Same input, same output & Yes & Not guaranteed \\
Behavior fully testable & Yes & Only probabilistically \\
Failure root cause & Clear & Often unclear \\
Compliance verification & Rule inspection & Behavior monitoring \\
\bottomrule
\end{tabular}
\end{table}

\subsection{The Fundamental Trust Challenge}
\label{subsec:fundamental-challenge}

This probabilistic nature creates a trust paradox:

\begin{quote}
\textbf{The more capable an agent becomes, the less predictable its behavior, and the harder it is to trust.}
\end{quote}

Sophisticated reasoning enables agents to handle novel situations---but also makes their behavior harder to anticipate. This is fundamentally different from traditional systems where capability and predictability are aligned.

\subsection{Empirical Evidence: LLM Reasoning Failures}
\label{subsec:empirical-evidence}

This theoretical concern has empirical grounding. Riddell et al. \citep{riddell2026} systematically evaluated six LLMs across 48,000 simulated cloud root cause analysis scenarios and developed a taxonomy of 16 distinct reasoning failure modes:

\begin{itemize}
    \item \textbf{Stalled:} Agent enters reasoning loops, unable to progress
    \item \textbf{Biased:} Agent anchors on irrelevant or misleading signals
    \item \textbf{Confused:} Multi-hop reasoning chains break down under complexity
\end{itemize}

The critical finding: these failures are \textit{structurally unpredictable}. Models that succeed on one scenario fail on structurally similar ones. Intermediate reasoning errors propagate to incorrect final diagnoses. This empirical work validates the theoretical framing: trust cannot be assumed based on capability demonstrations alone.

\section{Summary}
\label{sec:trust-problem-summary}

The fundamental challenge: agents are probabilistic, not deterministic. Traditional automation trust models assume predictability that agents cannot provide. This necessitates a new framework for reasoning about appropriate trust levels---developed in the next chapter.
